\section{序論}
11月9日朝の観測窓では、主要フロンティアラボによる新モデルの公式発表よりも、既存研究の解釈と次期モデルの公開時期を巡る予測が目立った。METRによるAnthropicの研究開発加速評価と、NVIDIAのCosmos 3 Nanoに関する議論は、AIが研究・物理シミュレーションの現場へ広がる動きを示している。

一方で、Sunoを巡る音声・著作権訴訟、データセンターの電力確保、ローカルコーディングエージェントの更新など、技術の運用条件にも関心が集まった。以下はGrokの収集順に記事化したもので、一次資料の独立検証は行っていない。
